\section{Alternative midpoint estimates for the segment forest}
\label{met:forest-variants}
The joint quadratic estimate gives the strongest forest bound in this
paper. The following inputs come from different intrinsic arguments and
illustrate the transfer rules of Section~\ref{met:transfer-section}.
They control radial deficit, parameterized test gains, or tails selected
by sequential compactness and a common primal norming test.

\subsection{Radial deficit}
\begin{corollary}[The atomic forest bound]\label{met:atomic-distortion}
For the segment forest space $X$, every map satisfying~\eqref{met:all-pairs}
has $D^2\ge1+k/16$.
\end{corollary}
\begin{proof}
For the forest atomic estimate, the constant
$C_*=\sqrt2+\sqrt6<4$ in the companion's arbitrary-displacement radial estimate
\cite[Proposition~12.1]{SegmentLenses} gives
\eqref{met:radial-input}. In particular,
\[
 \varepsilon^2/64\le R(R-\|x\|)\le R^2-\|x\|^2,
 \qquad D^2\ge1+k/16.
\]
This radial input gives a separate proof from the joint quadratic estimate.
\end{proof}
Keeping the
exact constant in the same calculation also gives
$k\le C_*^2D(D-1)$. Different atomic derivations with their own radial
constant $C$ give \eqref{met:additive-bound} with that constant; replacing
one proof's $C$ by another proof's value would discard its quantitative
content.

\subsection{Parameterized clipping tests}
Apply Proposition~\ref{met:parametric} with $a=1/2$ and $b=3$, as in the clipping input of
\cite[Section~10.1, Equation~(10.9)]{SegmentLenses}. Take $\lambda=1/(6D)$;
$Q-1\le3\lambda^2/2$ gives $\Delta\ge1/(24D)$.
Thus the full parameterized consequence is
\begin{equation}\label{met:parametric-twentyfour}
 1+\frac{k}{24D}\le D.
\end{equation}
For an optimized input
$\max_\pm\|p\pm y\|\ge\sqrt{\|p\|^2+a\|y\|^2}$,
the same limiting selection gives a child supremum at least
$\sqrt{\|x\|^2+a}$. Since
$\sqrt{r^2+a}-r$ decreases with $r>0$, some actual child has norm
at least $\|x\|+g/2$, where $g=\sqrt{D^2+a}-D>0$.
With $a=1/12$, as in
\cite[Proposition~10.4 and Equation~(10.10)]{SegmentLenses}, this gives
$1+kg/2\le D$. For this value of $a$, one
may make the finite selection before passing to a limit by taking approximation error
$\eta=g/8$: the norm comparison costs at most $4\eta$, leaving
$g/2$. Both descriptions retain the same actual-child loss.

Likewise, a radial input with constant $C$ gives a child supremum at
least $\|x\|+1/(C^2D)$. Selecting an actual child above the smaller
threshold $\|x\|+1/(2C^2D)$ gives
$1+k/(2C^2D)<D$ for $k\ge1$. The forest estimate in
\cite[Proposition~20.1]{SegmentLenses} permits $C=8$, giving the distinct
supremal gain $1/(64D)$ and actual gain greater than $1/(128D)$.

\subsection{An additive gain for a fixed distortion bound}
\label{met:bridge:fixed-bound-section}
The two-test estimate for an infinite bounded separated family in the segment
forest has the following form \cite[Proposition~11.3]{SegmentLenses}: for each $D\ge1$
there is $c_D>0$ such that
\begin{equation}\label{met:bridge:fixed-bound-input}
 \begin{gathered}
 1\le\|x\|\le D,\qquad \|z_i\|\le2D,\qquad
 \|z_i-z_j\|\ge1\quad(i\ne j)\\
 \Longrightarrow\quad
 \max_\pm\|x\pm z_i\|>\|x\|+c_D
 \quad\text{for some }i.
 \end{gathered}
\end{equation}
For fixed $D$, the same $c_D$ works for every center and every family
satisfying the displayed assumptions. The selected deviation need not
annihilate the head.

\begin{proposition}\label{met:bridge:fixed-bound}
If a target satisfies \eqref{met:bridge:fixed-bound-input}, every
noncontracting embedding of $\mathcal D_k$ with upper factor $D$ obeys
\[
 1+kc_D\le D.
\]
\end{proposition}
\begin{proof}
For every parent slope $x$ and its deviations $z_i$ from
\eqref{met:slopes}, one has $1\le\|x\|\le D$,
$\|z_i\|\le D\le2D$, and $\|z_i-z_j\|\ge2\ge1$.
The input therefore selects an actual child slope with norm greater
than $\|x\|+c_D$. Starting at the original slope and making this
choice through $k$ replacements gives a final slope greater than
$1+kc_D$ when $k\ge1$, and every final slope is at most $D$.
For $k=0$ the conclusion is $1\le D$.
\end{proof}

\subsection{Coordinate-null and sequential estimates}
Use the upper-one normalization with $\alpha=1/D$. At a parent with
child bound $0<R\le1$, its deviations satisfy $\|x\pm z_i\|\le R$
and $\|z_i-z_j\|\ge2\alpha$ for $i\ne j$.
On the countable forest coordinate set they admit a subsequence whose coordinates
all converge, by successive extraction and a diagonal choice. Pairing
successive terms gives $w_j=(z_{2j}-z_{2j-1})/2$ with coordinates tending
to zero, $\|w_j\|\ge\alpha$, and $\|x\pm w_j\|\le R$.
If the intrinsic estimate gives
\[
 \limsup_j\|w_j\|\le C\sqrt{R(R-\|x\|)},
\]
then
\[
 \|x\|\le R\left(1-\frac{\alpha^2}{C^2R^2}\right)
 \le R(1-\alpha^2/C^2),\qquad0<R\le1.
\]
The forest coordinate-null estimate
\cite[Corollary~23.5]{SegmentLenses} has
$C=\sqrt2+2\sqrt6<7$. Using $C=7$ retains the
$R$-dependent intermediate inequality and gives the backward factor
$1-1/(49D^2)$. Coordinatewise convergence here is used only with the
stated intrinsic estimate; it is not silently replaced by weak convergence.


The qualitative sequential cancellation theorem for the forest
\cite[Lemma~23.2]{SegmentLenses} satisfies the varying-head hypothesis
of Lemma~\ref{met:bridge:sequential-lens}. Hence it also gives, for each
fixed $D$, a uniform unit-lens gap and a backward bound $D^{-1}\le\rho^k$.
This application uses only sequential vanishing, without a quantitative
power law.

\subsection{Choosing the norming test before the head and the pair}
\label{met:bridge:normer-section}
For the segment forest, the weighted endpoint estimate in the
companion \cite[Theorem~16.1 and equation~(16.2)]{SegmentLenses} gives the following particular input.
Let $f$ be a finitely supported primal vector with $\rho(f)\le1$,
let $H$ be finite and ancestral with $\operatorname{supp}f\subset H$,
and suppose $z_+,z_-\in B_X$. Put
\[
 x=(z_++z_-)/2,\qquad
 d=1-x(f),\qquad a=\|Q_Hx\|.
\]
Then $d\ge0$ and
\begin{equation}\label{met:bridge:weighted-input}
 \|Q_Hz_\pm\|
 \le2\sqrt{(a+\sqrt{2d})^2+4d}+2\sqrt d.
\end{equation}
The same test $f$ is used for both endpoint laws; their individual
deficits are nonnegative and at most $2d$.

\begin{proposition}\label{met:bridge:normer-first}
For every $\varepsilon>0$ there is $0<\rho<1$ such that an infinite
$\varepsilon$-separated symmetric unit lens in the segment forest
space has center norm at most $\rho$.
\end{proposition}
\begin{proof}
Otherwise choose centers $x_n\in B_X$ with $\|x_n\|\to1$ and
infinite $\varepsilon$-separated families $(v_{n,i})_i$ satisfying
$\|x_n\pm v_{n,i}\|\le1$. First choose a finitely supported primal
near-normer $f_n$ with
\[
 \rho(f_n)\le1,\qquad x_n(f_n)>\|x_n\|-1/n.
\]
Next choose a finite ancestral $H_n$ containing
$\operatorname{supp}f_n$ and satisfying $\|Q_{H_n}x_n\|<1/n$.
This is possible by density of finite coordinate vectors in $X$ and
contractivity of $Q_{H_n}$: include the support of an approximation
to $x_n$ before taking its finite ancestral closure.

Only now choose two members of the $n$th lens whose projections
onto $H_n$ differ by less than $1/n$, and set
$u_n=(v_{n,i}-v_{n,j})/2$. Finite-dimensional total boundedness
justifies this choice, and crossed convex averages give
\begin{equation}\label{met:bridge:normer-selection}
 \|x_n\pm u_n\|\le1,\qquad
 \|P_{H_n}u_n\|<\frac1{2n},\qquad
 \|Q_{H_n}u_n\|\ge\frac\varepsilon2-\frac1{2n}.
\end{equation}
Apply \eqref{met:bridge:weighted-input} with the already chosen
$f_n,H_n$ and with $z_\pm=x_n\pm u_n$. Here
$d_n=1-x_n(f_n)\to0$ and $a_n=\|Q_{H_n}x_n\|\to0$.
Consequently both endpoint tails tend to zero. Their half-difference
$Q_{H_n}u_n$ also tends to zero, contradicting
\eqref{met:bridge:normer-selection}. No head has been chosen from
the selected pair; the order is near-normer, head, and then pair.
\end{proof}

For the original upper-one graph normalization, retain the separation
parameter $\varepsilon=1/D$. The normalized sibling deviations have
separation at least $2/(DL)\ge1/D$ when the child bound is
$0<L\le1$. The proposition gives the parent bound $\rho L$.
Backward iteration therefore gives $1/D\le\rho^k$, with one fixed
$\rho$ for this $D$.
